% ============================================================================
%  《Between Electric-Magnetic Duality and the Langlands Program》中译本
%  原书：David Ben-Zvi 教授 2021 年春季 UT Austin 课程讲义
%        笔记：Jackson Van Dyke
%  排版：XeLaTeX（请连续编译两遍以生成目录与交叉引用）
% ============================================================================
\documentclass[UTF8,11pt,a4paper,fontset=windows]{ctexbook}

% ---------- 页面与版式 ----------
\usepackage{geometry}
\geometry{left=2.55cm,right=2.55cm,top=2.55cm,bottom=2.6cm,headsep=0.5cm,footskip=1.1cm}
\usepackage{fancyhdr}
\usepackage{xcolor}
\usepackage{graphicx}
\usepackage{amsmath,amssymb}
\usepackage{amsthm}
\usepackage{caption}
\captionsetup{font=small,labelfont=bf,labelsep=quad}
\usepackage{booktabs}
\usepackage{array}
\usepackage{longtable}

% ---------- 超链接 ----------
\usepackage[unicode,bookmarksnumbered=true,breaklinks=true,
            colorlinks=true,linkcolor=blue!45!black,citecolor=blue!45!black,
            urlcolor=blue!45!black]{hyperref}
\hypersetup{
  pdftitle={电磁对偶与 Langlands 纲领之间（中译本·学术注疏版）},
  pdfauthor={David Ben-Zvi（讲义）；Jackson Van Dyke（笔记）；DeepSeek 学术助手（中译与注疏）},
  pdfsubject={几何 Langlands 纲领、电磁对偶、拓扑场论},
  pdfkeywords={Langlands program, electric-magnetic duality, TFT, gauge theory}
}

% ---------- 行距与段落 ----------
\linespread{1.18}
\setlength{\parskip}{0.25em}
\setlength{\parindent}{2em}
\setlength{\emergencystretch}{3em}

% ---------- 章节标题样式 ----------
\ctexset{
  chapter = {
    name = {第,章},
    number = \arabic{chapter},
    format = \Large\bfseries\centering,
    beforeskip = 8pt,
    afterskip = 18pt,
  },
  section = {
    format = \large\bfseries\raggedright,
    beforeskip = 14pt,
    afterskip = 6pt,
  },
  subsection = {
    format = \normalsize\bfseries\raggedright,
    beforeskip = 10pt,
    afterskip = 4pt,
  },
}

% ---------- 页眉页脚 ----------
\pagestyle{fancy}
\fancyhf{}
\fancyhead[LE]{\small\thepage}
\fancyhead[RO]{\small\thepage}
\fancyhead[LO]{\small\nouppercase{\leftmark}}
\fancyhead[RE]{\small\nouppercase{\rightmark}}
\renewcommand{\headrulewidth}{0.3pt}
\fancypagestyle{plain}{\fancyhf{}\fancyfoot[C]{\small\thepage}\renewcommand{\headrulewidth}{0pt}}

% ---------- 图片搜索路径 ----------
\graphicspath{{figures/}{figures/inline/}{./}}

% ---------- 术语宏 ----------
% 首次出现时使用：\term{中文译名}{English term}
\newcommand{\term}[2]{#1（#2）}
% 保留英文原词
\newcommand{\eng}[1]{\textit{#1}}

% ---------- 内嵌示意图 ----------
% 用法：\diagram[0.25\textwidth]{p135_d3.png}
\newcommand{\diagram}[2][0.28\textwidth]{%
  \begin{center}\includegraphics[width=#1]{#2}\end{center}}
% 带原文编号的图示方程：\diageq{10.13}{0.16\textwidth}{p135_d3.png}
\newcommand{\diageq}[3]{%
  \begin{equation}\tag{#1}\text{\includegraphics[width=#2]{#3}}\end{equation}}

% ---------- “本章重点” ----------
\newenvironment{keytakeaways}
  {\par\medskip\noindent
   \begin{minipage}{\linewidth}
   \hrule\vspace{5pt}\noindent\textbf{本章重点}\par\vspace{3pt}\small\noindent\ignorespaces}
  {\vspace{4pt}\hrule\end{minipage}\medskip}

% ---------- “本节要点” ----------
\newenvironment{sectionkey}
  {\par\smallskip\noindent
   \begin{minipage}{\linewidth}\small\noindent\textbf{本节要点：}\ignorespaces}
  {\end{minipage}\par\smallskip}

% ---------- 源文编号环境（保留原书编号） ----------
% 用法：\begin{srcdef}{3} ... \end{srcdef}
\newcommand{\srcenv}[3]{%
  \par\medskip\noindent\textbf{#1 #2.}\ \ignorespaces}
\newenvironment{srcdef}[1]{\par\medskip\noindent\textbf{定义 #1.}\ \ignorespaces}{\par\medskip}
\newenvironment{srcthm}[1]{\par\medskip\noindent\textbf{定理 #1.}\ \ignorespaces}{\par\medskip}
\newenvironment{srclem}[1]{\par\medskip\noindent\textbf{引理 #1.}\ \ignorespaces}{\par\medskip}
\newenvironment{srcprop}[1]{\par\medskip\noindent\textbf{命题 #1.}\ \ignorespaces}{\par\medskip}
\newenvironment{srccor}[1]{\par\medskip\noindent\textbf{推论 #1.}\ \ignorespaces}{\par\medskip}
\newenvironment{srcex}[1]{\par\medskip\noindent\textbf{例 #1.}\ \ignorespaces}{\par\medskip}
\newenvironment{srcrem}[1]{\par\medskip\noindent\textbf{注 #1.}\ \ignorespaces}{\par\medskip}
\newenvironment{srcconj}[1]{\par\medskip\noindent\textbf{猜想 #1.}\ \ignorespaces}{\par\medskip}
\newenvironment{srcclaim}[1]{\par\medskip\noindent\textbf{断言 #1.}\ \ignorespaces}{\par\medskip}
\newenvironment{srcnotation}[1]{\par\medskip\noindent\textbf{记号 #1.}\ \ignorespaces}{\par\medskip}

% ---------- 常用数学宏 ----------
\newcommand{\C}{\mathbb{C}}
\newcommand{\R}{\mathbb{R}}
\newcommand{\Z}{\mathbb{Z}}
\newcommand{\N}{\mathbb{N}}
\newcommand{\Q}{\mathbb{Q}}
\newcommand{\F}{\mathbb{F}}
\newcommand{\HH}{\mathbb{H}}
\newcommand{\GG}{\mathbb{G}}
\newcommand{\PP}{\mathbb{P}}
\newcommand{\bA}{\mathbb{A}}
\newcommand{\cO}{\mathcal{O}}
\newcommand{\cF}{\mathcal{F}}
\newcommand{\cG}{\mathcal{G}}
\newcommand{\cH}{\mathcal{H}}
\newcommand{\cC}{\mathcal{C}}
\newcommand{\cD}{\mathcal{D}}
\newcommand{\cA}{\mathcal{A}}
\newcommand{\cB}{\mathcal{B}}
\newcommand{\cL}{\mathcal{L}}
\newcommand{\cM}{\mathcal{M}}
\newcommand{\cN}{\mathcal{N}}
\newcommand{\cX}{\mathcal{X}}
\newcommand{\cZ}{\mathcal{Z}}
\newcommand{\fg}{\mathfrak{g}}
\renewcommand{\proofname}{证明}
\newcommand{\sph}{\mathrm{Sph}}
\newcommand{\Bun}{\mathrm{Bun}}
\newcommand{\Loc}{\mathrm{Loc}}
\newcommand{\Rep}{\mathrm{Rep}}
\newcommand{\Vect}{\mathrm{Vect}}
\newcommand{\Shv}{\mathrm{Shv}}
\newcommand{\End}{\mathrm{End}}
\newcommand{\Hom}{\mathrm{Hom}}
\newcommand{\Aut}{\mathrm{Aut}}
\newcommand{\Spec}{\mathrm{Spec}}
\newcommand{\Gal}{\mathrm{Gal}}
\newcommand{\Gr}{\mathrm{Gr}}
\newcommand{\QC}{\mathrm{QC}}
\newcommand{\IndCoh}{\mathrm{IndCoh}}
\newcommand{\Bord}{\mathrm{Bord}}
\newcommand{\Pic}{\mathrm{Pic}}
\newcommand{\Jac}{\mathrm{Jac}}
\newcommand{\G}{\mathbb{G}}
\newcommand{\cE}{\mathcal{E}}
\newcommand{\cP}{\mathcal{P}}
\newcommand{\Sph}{\mathrm{Sph}}
\newcommand{\GL}{\mathrm{GL}}
\newcommand{\SL}{\mathrm{SL}}
\newcommand{\PSL}{\mathrm{PSL}}
\newcommand{\PGL}{\mathrm{PGL}}
\newcommand{\Sp}{\mathrm{Sp}}
\newcommand{\SO}{\mathrm{SO}}
\newcommand{\LG}{\mathrm{LG}}

% ---------- 图表标题前缀 ----------
\renewcommand{\figurename}{图}
\renewcommand{\tablename}{表}
\renewcommand{\contentsname}{目\quad 录}
\renewcommand{\bibname}{参考文献}

% ---------- 微调 ----------
\allowdisplaybreaks[2]
\setlength{\abovedisplayskip}{6pt}
\setlength{\belowdisplayskip}{6pt}
